% TEMPLATE-MILC-v3.tex - plantilla maestra UNICA de las guias MILC v3.
%
% La usan las tres familias de guias, cada una con su driver:
%   · Grados 6-11  -> scripts/build-guias-g11.py
%   · Bebras       -> scripts/build-guias-bebras.py
%   · Semillero    -> scripts/build-guias-semillero.py
%
% Antes habia tres copias casi identicas (~400 lineas iguales, 6 distintas):
% cada mejora habia que aplicarla tres veces, y en la practica no se hacia
% (el motor de recursos vivia solo en dos de ellas). Lo que varia entre
% familias esta parametrizado en 6 placeholders que cada driver llena
% (se nombran aqui SIN su sintaxis real: el builder reemplaza por texto plano,
% tambien dentro de los comentarios, y un valor multilinea rompe el preambulo):
%
%   HEADER_IZQ        encabezado izquierdo de pagina
%   HEADER_DER        encabezado derecho de pagina
%   PORTADA_ETIQUETA  rotulo sobre el numero grande (GRADO/MOMENTO/MODULO)
%   PORTADA_SERIE     linea de serie de la portada
%   PORTADA_ID        identificador de la guia en la portada
%   PORTADA_CIRCULO   el circulito de grado (°), o vacio si no aplica
%
% El resto de claves las documenta content/guias/_SCHEMA.md. El script valida
% que no queden placeholders sin reemplazar antes de escribir el archivo final.

\documentclass[11pt,letterpaper]{article}
\usepackage[margin=2.0cm]{geometry}
\usepackage[table]{xcolor}
\usepackage{fontspec}
\usepackage[spanish]{babel}
\usepackage{tikz}
% Librerías TikZ para los diagramas de las guías (bloque `recursos:`):
%  · babel  → babel en español vuelve activos caracteres como «>», y TikZ los usa
%             en sus opciones (>=stealth, ->). Sin esta librería, cualquier
%             diagrama con flechas revienta con «\language@active@arg> has an
%             extra }». Debe ir ANTES de que se dibuje nada.
%  · positioning        → sintaxis «below left=2pt and 3pt».
%  · decorations.pathreplacing → llaves ({brace}) para señalar rangos.
\usetikzlibrary{babel,positioning,decorations.pathreplacing}
\usepackage{graphicx}
% Circuitos: el trabajo de electrónica se hace hoy con micro:bit y sus sensores
% integrados (MakeCode, sin cableado), así que no se necesitan esquemáticos.
% Si algún día entran montajes con protoboard, basta descomentar esta línea:
% los diagramas .tex/.tikz declarados en `recursos:` ya se incrustan solos.
% \usepackage{circuitikz}
\usepackage[most]{tcolorbox}
\usepackage{tabularx}
\usepackage{array}
\usepackage{fancyhdr}
\usepackage{titlesec}
\usepackage[document]{ragged2e}  % bandera a la derecha en todo el cuerpo
\usepackage{enumitem}
\usepackage{microtype}

% Helvetica Neue no trae la flecha U+2192, asi que cada «→» escrito literalmente
% en el YAML se compilaba a NADA, en silencio (462 flechas en 61 guias: rutas del
% tipo «paleta → tipografia → lockup» salian rotas). Mapeamos el caracter a su
% version matematica, que si tiene glifo. Asi el autor puede seguir escribiendo
% «→» comodamente en el YAML.
\usepackage{newunicodechar}
\newunicodechar{→}{\ensuremath{\rightarrow}}
% Mismo problema con los simbolos de marca y vinetas: el «✓» de «una palomita ✓»
% salia como cuadrito vacio en el PDF de todas las guias que lo usaban.
\usepackage{pifont}
\newunicodechar{✓}{\ding{51}}
\newunicodechar{✗}{\ding{55}}
\newunicodechar{✔}{\ding{52}}
\newunicodechar{•}{\textbullet}
\newunicodechar{≥}{\ensuremath{\geq}}
\newunicodechar{≤}{\ensuremath{\leq}}

\setmainfont{Helvetica Neue}
\setsansfont{Helvetica Neue}
\setlength{\parindent}{0pt}
\setlength{\parskip}{6pt}
% Sin particion de palabras: en 6.º-7.º una palabra cortada por guion al final
% de linea («Comuni-cacion») frena la lectura mucho mas que una linea corta.
\hyphenpenalty=10000
\exhyphenpenalty=10000
\pretolerance=2000
\tolerance=3000
\setlength{\emergencystretch}{3em}
\linespread{1.10}
\renewcommand{\arraystretch}{1.25}
\setlist{leftmargin=5mm,itemsep=1mm,topsep=1mm}
\newcolumntype{Y}{>{\RaggedRight\arraybackslash}X}

\definecolor{milcMagenta}{HTML}{E6007E}
\definecolor{milcVino}{HTML}{5A0038}
\definecolor{milcTurquesa}{HTML}{0093A5}
\definecolor{milcVerde}{HTML}{58A923}
\definecolor{milcMostaza}{HTML}{E5B400}
\definecolor{milcOcre}{HTML}{B86600}
\definecolor{milcNegro}{HTML}{241816}
\definecolor{milcGris}{HTML}{F7F7F4}
\definecolor{milcLinea}{HTML}{E8E2DD}
\definecolor{milcGradePrimary}{HTML}{0066FF}
\definecolor{milcGradeDark}{HTML}{003D99}
\definecolor{milcGradeSoft}{HTML}{D6E8FF}
\definecolor{milcGradeAccent}{HTML}{CFE7FF}
\definecolor{milcGradeText}{HTML}{FFFFFF}
\color{milcNegro}

\pagestyle{fancy}
\fancyhf{}
\lhead{\small Tecnología e Informática · Grado 6}
\rhead{\small Guía 20 · MILC v3}
\cfoot{\small\thepage}
\renewcommand{\headrulewidth}{0pt}

\newtcolorbox{titlebox}[1]{
  enhanced,colback=#1,colframe=#1,arc=7mm,boxrule=0pt,
  left=7mm,right=7mm,top=3mm,bottom=3mm,
  before skip=8pt,after skip=8pt,
  fontupper=\sffamily\bfseries\Large\color{white}
}

\newtcolorbox{softbox}[3]{
  enhanced,colback=#2,colframe=#1,borderline west={4pt}{0pt}{#1},
  arc=2mm,boxrule=.4pt,left=4mm,right=4mm,top=3mm,bottom=3mm,
  before skip=6pt,after skip=8pt,title={#3},coltitle=white,
  colbacktitle=#1,fonttitle=\sffamily\bfseries,fontupper=\RaggedRight,
  attach boxed title to top left={xshift=3mm,yshift=-2mm},
  boxed title style={arc=2mm, boxrule=0pt}
}

\newtcolorbox{drawbox}[2]{
  enhanced,colback=white,colframe=#1,arc=4mm,boxrule=1pt,
  left=5mm,right=5mm,top=4mm,bottom=4mm,before skip=8pt,after skip=8pt,
  title={#2},coltitle=white,colbacktitle=#1,fonttitle=\sffamily\bfseries,
  fontupper=\RaggedRight,
  attach boxed title to top left={xshift=4mm,yshift=-2mm},
  boxed title style={arc=3mm, boxrule=0pt}
}

\newtcolorbox{infoband}[1]{
  enhanced,colback=#1!8,colframe=#1!50,arc=2mm,boxrule=.5pt,
  left=4mm,right=4mm,top=2mm,bottom=2mm,before skip=4pt,after skip=4pt,
  fontupper=\small\RaggedRight
}

% Puente narrativo entre secciones (contrato editorial Conéctate).
% Da continuidad: "Acabas de X. Ahora vas a Y porque Z."
% Visualmente: banda izquierda en color del grado, texto en itálica pequeña.
\newtcolorbox{puentebox}{
  enhanced,colback=milcGradeSoft,colframe=milcGradeDark,
  borderline west={3pt}{0pt}{milcGradeDark},
  arc=0mm,boxrule=0pt,left=5mm,right=4mm,top=2.5mm,bottom=2.5mm,
  before skip=4pt,after skip=8pt,
  fontupper=\itshape\small\color{milcNegro}\RaggedRight
}
% La flecha va en modo matematico a proposito: Helvetica Neue NO tiene el glifo
% U+2192, asi que \textrightarrow se compilaba a NADA («Missing character: There
% is no → in font Helvetica Neue») y el puente salia sin flecha. En modo
% matematico la flecha viene de la fuente matematica y si se dibuja.
\newcommand{\puente}[1]{%
  \begin{puentebox}%
    \textcolor{milcGradeDark}{$\rightarrow$}\hspace{2mm}#1%
  \end{puentebox}%
}

\newcommand{\linead}[1]{\noindent\color{milcLinea}\rule{#1}{0.5pt}\color{milcNegro}}
\newcommand{\respuesta}[1]{\par\vspace{2mm}\foreach \n in {1,...,#1}{\noindent\color{milcLinea}\rule{\linewidth}{0.5pt}\par\vspace{5mm}}\color{milcNegro}}
\newcommand{\checkbox}{\textcolor{milcGradeDark}{\fbox{\phantom{x}}}\hspace{5pt}}

\newcommand{\phasecard}[4]{
\begin{tcolorbox}[enhanced,colback=#3,colframe=#2,arc=3mm,boxrule=.5pt,height=3.05cm,valign=center,halign=center,left=1.5mm,right=1.5mm,top=2mm,bottom=2mm]
{\sffamily\bfseries\fontsize{22}{24}\selectfont\textcolor{#2}{#1}}\par
{\sffamily\bfseries\small #4}
\end{tcolorbox}
}

% ── Piezas del rediseno 2026 (legibilidad 6.º-11.º) ───────────────────
% Banda de estacion: reemplaza el titlebox macizo. Titulo a la izquierda,
% dato practico (tiempo/modalidad) a la derecha, en una sola linea.
\newcommand{\milcestacion}[2]{%
\begin{tcolorbox}[enhanced,colback=#1,colframe=#1,arc=2mm,boxrule=0pt,
  left=5mm,right=5mm,top=2.6mm,bottom=2.6mm,before skip=11pt,after skip=7pt]
{\sffamily\bfseries\fontsize{15}{18}\selectfont\color{white}#2}
\end{tcolorbox}}

% Tarjeta de paso numerada: sustituye las tablas «Pilar / Aplicacion», que
% eran formato de consulta docente, no de estudiante. El numero va en un
% circulo sobre el borde para que la secuencia se lea de un vistazo.
\newcommand{\milcpaso}[3]{%
\begin{tcolorbox}[enhanced,colback=white,colframe=milcLinea,arc=1.5mm,boxrule=.9pt,
  left=14mm,right=4mm,top=3mm,bottom=3mm,before skip=4pt,after skip=4pt,
  fontupper=\RaggedRight,
  overlay={\node[anchor=north west,circle,fill=#2,text=white,inner sep=0pt,
    minimum size=7.4mm,font=\sffamily\bfseries\small]
    at ([xshift=3.6mm,yshift=-3.2mm]frame.north west) {#1};}]
#3
\end{tcolorbox}}

% Casilla marcable de tamano util con lapiz (la anterior era \fbox{\phantom{x}},
% de alto variable segun el texto que la seguia).
\newcommand{\milcmarca}{\framebox[3.8mm]{\rule{0pt}{3.4mm}}}

% Fila marcable de lista suelta (checks de la estacion 1).
\newcommand{\milccheckrow}[1]{%
\par\vspace{1.8mm}\noindent
\makebox[6.4mm][l]{\raisebox{-.55mm}{\textcolor{milcNegro}{\milcmarca}}}%
\parbox[t]{\dimexpr\linewidth-6.4mm\relax}{\RaggedRight #1}\par}

% Celda marcable dentro de la rubrica (tabla de autoevaluacion). Misma
% sangria francesa, pero pensada para vivir dentro de una columna X.
\newcommand{\milcopcion}[1]{%
\makebox[5.2mm][l]{\raisebox{-.55mm}{\textcolor{milcNegro}{\milcmarca}}}%
\parbox[t]{\dimexpr\linewidth-5.2mm\relax}{\RaggedRight #1}}

% ── Recursos visuales de la guía ──────────────────────────────────────
% Declarados en el bloque `recursos:` del YAML y emitidos por el builder.
% \guiaFigura   → imagen rasterizada o PDF (foto, carta celeste, montaje).
% \guiaDiagrama → fuente TikZ/circuitikz (.tex) incrustada como vector:
%                 es la vía para circuitos y esquemáticos editables.
% #1 = ruta relativa al .tex · #2 = pie (puede ir vacío)
\newcommand{\guiaFigura}[2]{%
\begin{center}
\includegraphics[width=0.88\linewidth,height=0.38\textheight,keepaspectratio]{#1}\\[1.5mm]
{\footnotesize\itshape\textcolor{milcNegro!75}{#2}}
\end{center}
}

\newcommand{\guiaDiagrama}[2]{%
\begin{center}
\input{#1}\\[1.5mm]
{\footnotesize\itshape\textcolor{milcNegro!75}{#2}}
\end{center}
}

\begin{document}
\thispagestyle{empty}

% ── Portada ───────────────────────────────────────────────────────────
% Antes: un rectangulo de color a pagina completa. Se veia bien en pantalla,
% pero en la fotocopiadora del colegio se comia el toner y salia gris sucio.
% Ahora la banda ocupa el tercio superior y el resto va en blanco.
\begin{tikzpicture}[remember picture,overlay]
  \path[fill=milcGradePrimary]
    (current page.north west) rectangle ([yshift=-10.6cm]current page.north east);

  \node[anchor=north west,text=milcGradeText] at ([xshift=2.0cm,yshift=-1.55cm]current page.north west)
    {\sffamily\bfseries\fontsize{12}{14}\selectfont GRADO};
  \node[anchor=north west,text=milcGradeText,opacity=.85] at ([xshift=2.0cm,yshift=-2.35cm]current page.north west)
    {\sffamily\bfseries\fontsize{8.5}{10}\selectfont SERIE GUÍAS · TECNOLOGÍA E INFORMÁTICA};
  \node[anchor=north east,text=milcGradeText,opacity=.85,align=right] at ([xshift=-2.0cm,yshift=-1.55cm]current page.north east)
    {\sffamily\bfseries\fontsize{9}{12}\selectfont I.E. Sor María Juliana\\Cartago, Valle del Cauca};

  \node[anchor=west,text=milcGradeText] (gradonum) at ([xshift=1.85cm,yshift=-7.1cm]current page.north west)
    {\sffamily\bfseries\fontsize{112}{112}\selectfont 6};
  % El circulito de grado (°) solo aplica a las guias de grado («11°»); en
  % Bebras y Semillero el numero grande es un momento/modulo, no un grado.
  % Se posiciona relativo al nodo `gradonum`, no en coordenadas absolutas.
  \draw[milcGradeText,line width=1.6pt]
    ([xshift=2.0mm,yshift=-4.5mm]gradonum.north east) circle (.15cm);

  \node[anchor=west,text=milcGradeText,text width=11.4cm,align=left] at ([xshift=1.5cm,yshift=0.5cm]gradonum.east)
    {
     {\sffamily\bfseries\fontsize{9}{11}\selectfont\MakeUppercase{Período 2 · Hardware y software}}\\[3mm]
     {\sffamily\bfseries\fontsize{21}{25}\selectfont\setlength{\baselineskip}{27pt}Mi computador\\[4pt]ideal ---\\[4pt]diseñar la\\[4pt]máquina que\\[4pt]yo necesito}\\[3.5mm]
     {\sffamily\bfseries\fontsize{15}{18}\selectfont Guía 20-6-TIC}
    };
\end{tikzpicture}

\vspace*{9.4cm}
\vfill

{\sffamily\bfseries\fontsize{8}{10}\selectfont\textcolor{milcGradeDark}{LO QUE TE LLEVAS HOY}}

\begin{tcolorbox}[enhanced,colback=white,colframe=milcNegro,arc=2mm,boxrule=1.6pt,
  left=5mm,right=5mm,top=4mm,bottom=4mm,before skip=3pt,after skip=12pt,
  fontupper=\RaggedRight\fontsize{11}{15.5}\selectfont]
Una \textbf{ficha técnica de tu computador ideal} en una hoja del cuaderno, con: \textbf{(a) perfil de uso} (¿para qué lo necesitas? colegio, juegos, edición, ofimática), \textbf{(b) lista de componentes recomendados} (CPU, RAM, disco, GPU, monitor, periféricos), \textbf{(c) presupuesto aproximado} (puedes consultar precios en internet o estimar), \textbf{(d) justificación con criterio} de por qué escogiste cada componente, y \textbf{(e) sugerencia de software inicial} (sistema operativo + 5 programas que vas a usar). Cierre: \emph{``Si tuviera el dinero, este es el equipo que compraría, y sé por qué.''}.
\end{tcolorbox}

\begin{tcolorbox}[enhanced,colback=milcGris,colframe=milcGris,arc=2mm,boxrule=0pt,
  left=5mm,right=5mm,top=3.5mm,bottom=3.5mm,after skip=12pt]
{\sffamily\bfseries\fontsize{8}{10}\selectfont\textcolor{milcNegro!55}{ESTA GUÍA ES DE}}\\[3mm]
\begin{tabularx}{\linewidth}{X p{3.2cm} p{3.2cm}}
\linead{\linewidth} & \linead{3.0cm} & \linead{3.0cm}\\[-1mm]
{\fontsize{8.5}{10}\selectfont\textcolor{milcNegro!55}{Nombre completo}} &
{\fontsize{8.5}{10}\selectfont\textcolor{milcNegro!55}{Grupo}} &
{\fontsize{8.5}{10}\selectfont\textcolor{milcNegro!55}{Fecha}}\\
\end{tabularx}
\end{tcolorbox}

\vfill

\noindent\textcolor{milcLinea}{\rule{\linewidth}{1pt}}\\[2mm]
{\sffamily\bfseries\fontsize{9.5}{12}\selectfont Autor: Dr. Álvaro Cárdenas Orozco}\\[1mm]
{\sffamily\fontsize{8.5}{11}\selectfont\textcolor{milcNegro!55}{Metodología MILC · apertura ancestral · pensamiento computacional · triángulo Dussel–estoicismo–Floridi}}

\newpage

\section*{Apertura: Mi computador ideal --- diseñar la máquina que yo necesito}

\begin{softbox}{milcGradeDark}{milcGradeSoft}{Saber ancestral}
\textbf{Don Lucho el relojero no vendía un solo modelo de reloj a todos sus clientes.} Cuando alguien entraba a la vitrina de la calle 14 de Cartago, lo primero que don Lucho hacía era \textbf{preguntar}: \emph{``¿Para qué lo quiere usted, mijo? ¿Para el trabajo? ¿Para regalar? ¿Para el campo o la ciudad?''}. Si era para un campesino que iba al cafetal, recomendaba uno \emph{resistente al sudor y al polvo}, sencillo. Si era para una secretaria, uno \emph{elegante de pulsera fina}. Si era para un abuelo que ya no veía bien, uno con \emph{números grandes}. \textbf{El relojero sabio sabe que no hay UN reloj para todos, hay EL reloj para cada quién}. Por eso preguntaba primero, vendía después. Hoy con los computadores pasa lo mismo. No hay UN computador mejor: hay \textbf{el computador correcto para lo que TÚ vas a hacer}. Diseñar tu equipo ideal es la prueba de fuego de todo lo aprendido en este periodo: combinar hardware, software, mantenimiento y ergonomía para decidir con criterio.
\end{softbox}

\begin{softbox}{milcTurquesa}{milcGris}{Saber contemporáneo}
\textbf{Diseñar el computador ideal} no es elegir el más caro --- es elegir el que se adapta a TU uso. Hay \textbf{4 perfiles típicos} y cada uno requiere componentes distintos. \textbf{(1) Ofimática y colegio}: navegar, escribir documentos, ver videos educativos. Computador básico, RAM 8 GB, disco SSD 256 GB, sin tarjeta gráfica dedicada. \textbf{(2) Edición media}: editar fotos, videos cortos, presentaciones complejas. RAM 16 GB, SSD 512 GB, tarjeta gráfica básica. \textbf{(3) Gamer}: jugar juegos pesados (Fortnite, Minecraft con shaders, GTA V). RAM 16-32 GB, SSD 1 TB, tarjeta gráfica dedicada (GTX o RTX). \textbf{(4) Creador profesional}: editar videos largos, programación, modelado 3D. RAM 32 GB+, SSD 1 TB+, tarjeta gráfica de gama alta. \textbf{Saber tu perfil te ahorra dinero}: si solo vas a hacer tareas, no necesitas un computador gamer; si vas a jugar, no compres un básico. Esta clase te entrena en \emph{decidir con criterio} en vez de \emph{dejarte llevar por la publicidad}.
\end{softbox}

\begin{softbox}{milcMostaza}{milcGris}{Pregunta-puente}
\textit{Si tus papás te dijeran \emph{``escoge tu primer computador propio, tienes 2 millones de pesos''}, ¿\textbf{qué le preguntarías al vendedor} primero? ¿Sabrías qué componentes pedir?}
\end{softbox}

\begin{softbox}{milcVerde}{milcGris}{Saber y saber hacer}
Al terminar la clase: (1) podrás \textbf{identificar} tu propio perfil de uso; (2) sabrás \textbf{aplicar} criterios de hardware a los 4 perfiles; (3) podrás \textbf{evaluar} si una oferta de computador conviene o no; (4) habrás \textbf{creado} tu ficha técnica con justificación.
\end{softbox}



\small
\begin{tabularx}{\textwidth}{>{\bfseries}p{3.0cm}Y}
\rowcolor{milcGradePrimary!12}Autor & Dr. Álvaro Cárdenas Orozco\\
DBA & Identifica componentes y sistemas tecnológicos en su entorno (MEN, T\&I 6°)\\
Referentes & Saber ancestral del relojero · Sistemas como cadena de oficios · Diagnóstico paso a paso\\
Producto final & Una \textbf{ficha técnica de tu computador ideal} en una hoja del cuaderno, con: \textbf{(a) perfil de uso} (¿para qué lo necesitas? colegio, juegos, edición, ofimática), \textbf{(b) lista de componentes recomendados} (CPU, RAM, disco, GPU, monitor, periféricos), \textbf{(c) presupuesto aproximado} (puedes consultar precios en internet o estimar), \textbf{(d) justificación con criterio} de por qué escogiste cada componente, y \textbf{(e) sugerencia de software inicial} (sistema operativo + 5 programas que vas a usar). Cierre: \emph{``Si tuviera el dinero, este es el equipo que compraría, y sé por qué.''}.\\
\end{tabularx}
\normalsize

\puente{Hoy haces 4 cosas: identificas tu perfil de uso, aprendes los 4 perfiles típicos y sus componentes, diseñas tu ficha técnica, y firmas tu propuesta.}

\milcestacion{milcGradeDark}{Ruta MILC: así vamos a aprender}
\begin{tabularx}{\textwidth}{>{\centering\arraybackslash}X>{\centering\arraybackslash}X>{\centering\arraybackslash}X>{\centering\arraybackslash}X}
\phasecard{1}{milcVerde}{milcGris}{Escucha\\[-1mm]\small Observo, escucho y pregunto.} &
\phasecard{2}{milcTurquesa}{milcGris}{Entender\\[-1mm]\small Sistematizo y ordeno.} &
\phasecard{3}{milcMagenta}{milcGris}{Hacer\\[-1mm]\small Construyo y aplico.} &
\phasecard{4}{milcVino}{milcGris}{Cierre\\[-1mm]\small Evalúo y reflexiono.}
\end{tabularx}


\puente{Empezamos preguntándote: ¿para qué usarías TÚ el computador?}

\milcestacion{milcVerde}{Estación 1 de 3 · Escucha}

\begin{softbox}{milcVerde}{milcGris}{Escena para escuchar}
\textbf{Actividad 1 · IDENTIFICA --- ¿Cuál es TU perfil?} (10 min · individual). En el cuaderno, responde estas \textbf{5 preguntas} con sinceridad: \emph{(1) ¿Qué actividades haces más en el computador (escoge las 3 principales): hacer tareas, ver videos, jugar, editar fotos/videos, programar, otra?} \emph{(2) ¿Cuánto tiempo lo usas al día?} \emph{(3) ¿Necesitas que sea portátil para llevarlo al colegio?} \emph{(4) ¿Tu casa tiene buen internet o necesitas un equipo que aguante mucho local?} \emph{(5) ¿Te gusta jugar videojuegos pesados o solo casuales?}. Estas respuestas determinan tu perfil. Sé honesto: si no juegas mucho, no inventes que sí solo porque suena \emph{cool}.
\end{softbox}

{\sffamily\bfseries\fontsize{8}{10}\selectfont\textcolor{milcNegro!55}{MARCA CUANDO LO HAYAS HECHO}}
\milccheckrow{Respondí las 5 preguntas con honestidad.}
\milccheckrow{Mis respuestas reflejan mi uso real, no un uso imaginado.}
\milccheckrow{Identifiqué cuáles 3 actividades son las que más hago.}

\begin{infoband}{milcVerde}


\textbf{Lo que va al cuaderno (Actividad 1 · IDENTIFICA).} Título: «Actividad 1 --- Mi perfil de usuario». \textbf{Formato:} 5 preguntas respondidas en bloques cortos. \textbf{Extensión:} 5 respuestas de 1-2 líneas cada una. \textbf{Sabes que terminaste cuando} tu perfil queda claro: si alguien lee esto, sabe para qué usas el computador.
\end{infoband}

\puente{Después aprendes los 4 perfiles con sus componentes recomendados y rangos de precio.}

\milcestacion{milcTurquesa}{Estación 2 de 3 · Entender}

\begin{softbox}{milcTurquesa}{milcGris}{Lo que vas a entender}
Las 5 respuestas te ubican en uno de los \textbf{4 perfiles típicos}. Conocer tu perfil te ayuda a saber qué hardware necesitas (y qué no). \textbf{No hay perfil mejor}: hay perfil correcto para cada quién.
\end{softbox}

\milcpaso{1}{milcTurquesa}{\textbf{Perfil 1 --- Ofimática + colegio (el más común para 6° grado).} Usas el computador para tareas, Word, navegar, ver videos educativos, oír música. \textbf{Componentes ideales:} CPU intermedia (Intel Core i3 o i5; AMD Ryzen 3 o 5). RAM 8 GB (lo justo). Disco SSD 256 GB (rápido, suficiente). Sin tarjeta gráfica dedicada (la del CPU basta). Monitor 14-15 pulgadas (portátil) o 22 pulgadas (escritorio). \textbf{Rango de precio en Colombia 2026:} 1.500.000 a 2.500.000 pesos.}
\milcpaso{2}{milcTurquesa}{\textbf{Perfil 2 --- Edición media (presentaciones, foto, videos cortos).} Usas el equipo para tareas Y edición ligera: fotos en Canva, videos cortos en CapCut, presentaciones con efectos. \textbf{Componentes ideales:} CPU media-alta (i5 o Ryzen 5). RAM 16 GB. Disco SSD 512 GB. Tarjeta gráfica básica integrada o dedicada de entrada (GTX 1650 o equivalente). Monitor de buena resolución (Full HD). \textbf{Rango de precio:} 2.500.000 a 4.000.000 pesos.}
\milcpaso{3}{milcTurquesa}{\textbf{Perfil 3 --- Gamer (juegos pesados regulares).} Usas el equipo para jugar Fortnite, Roblox con shaders, Minecraft modeado, GTA V, juegos con gráficos exigentes. \textbf{Componentes ideales:} CPU alta (i5 o i7; Ryzen 5 o 7). RAM 16 GB mínimo, 32 GB ideal. Disco SSD 1 TB. \textbf{Tarjeta gráfica dedicada} (GTX 1660, RTX 3060 o superior). Monitor con buena tasa de refresco (75 Hz+). \textbf{Rango de precio:} 4.000.000 a 7.000.000 pesos.}
\milcpaso{4}{milcTurquesa}{\textbf{Perfil 4 --- Creador profesional (videos largos, programación, 3D).} Usas el equipo para edición de video largo en Premiere o DaVinci, programación, modelado 3D en Blender. \textbf{Componentes ideales:} CPU alta (i7 o i9; Ryzen 7 o 9). RAM 32 GB+. Disco SSD 1 TB + HDD 2 TB para almacenamiento. Tarjeta gráfica de gama alta (RTX 3070+ o equivalente). Monitor 27 pulgadas o doble monitor. \textbf{Rango de precio:} 7.000.000 a 12.000.000+ pesos. \emph{Este perfil rara vez se necesita en 6°}; es para adultos con trabajo específico.}

\begin{drawbox}{milcTurquesa}{4 perfiles + componentes clave}
\textbf{P1 · Ofimática + colegio:} i3/i5, 8 GB RAM, SSD 256 GB, sin GPU. \\[1mm]
\textbf{P2 · Edición media:} i5, 16 GB RAM, SSD 512 GB, GPU básica. \\[1mm]
\textbf{P3 · Gamer:} i5/i7, 16-32 GB RAM, SSD 1 TB, GPU dedicada. \\[1mm]
\textbf{P4 · Creador pro:} i7/i9, 32 GB+ RAM, SSD 1 TB+ HDD 2 TB, GPU alta.
\end{drawbox}

\begin{softbox}{milcMostaza}{milcGris}{Errores comunes y su corrección}
\textbf{Errores frecuentes al elegir computador.} (1) \emph{``Compro el más caro y ya''} --- Falso. El más caro suele ser más de lo que necesitas. Pagas por componentes que nunca usarás. (2) \emph{``Compro el más barato''} --- Falso al otro extremo. Si es muy básico, se vuelve lento en 6 meses y lo terminas cambiando. (3) \emph{``Solo importa la marca''} --- Falso. Lo que importa son los componentes (CPU, RAM, disco). Una marca \emph{cool} con malos componentes es peor que una marca menos famosa con buenos componentes. (4) \emph{``Me dejo llevar por la publicidad del vendedor''} --- Falso. El vendedor quiere ganar comisión. Tú quieres el equipo correcto para ti. Lleva tu perfil escrito y pide componentes específicos.
\end{softbox}

\begin{infoband}{milcTurquesa}


\textbf{Lo que va al cuaderno (Actividad 2 · EXPLICA).} Título: «Actividad 2 --- Los 4 perfiles con sus componentes». \textbf{Formato:} tabla de 4 filas, 5 columnas. Columnas: perfil, CPU, RAM, disco, GPU. \textbf{Extensión:} 4 filas. \textbf{Sabes que terminaste cuando} puedes decir en voz alta qué componentes necesita cada perfil.
\end{infoband}

\puente{Con todo claro, diseñas la ficha técnica de tu computador ideal con justificación.}

\milcestacion{milcMagenta}{Estación 3 de 3 · Hacer}

\begin{softbox}{milcMagenta}{milcGris}{Cómo vas a construir tu producto}
\textbf{Actividad 3 · CREA --- Tu ficha técnica del computador ideal} (32 min · individual). En una hoja completa del cuaderno vas a hacer la \textbf{ficha técnica de TU computador ideal}. Combina lo aprendido en todo el periodo: hardware, software, mantenimiento, ergonomía. Esta es la cosecha del periodo 2.
\end{softbox}

\milcpaso{1}{milcMagenta}{\textbf{Paso 1 --- Título y mi perfil.} Arriba escribe: \textbf{``Mi computador ideal --- ficha técnica''} y debajo: \emph{``[Tu nombre], grado 6[tu letra], año 2026''}. Subraya el título. Luego escribe en un cuadro: \textbf{``Mi perfil:''} y completa con tu perfil de la Actividad 1 (ofimática, edición media, gamer, o mezcla). Sé honesto: \emph{``ofimática 60\% + gamer casual 40\%''} es válido.}
\milcpaso{2}{milcMagenta}{\textbf{Paso 2 --- Lista de componentes.} En una tabla de \textbf{6 filas, 3 columnas}, lista los componentes principales. Columnas: \emph{componente}, \emph{lo que elijo}, \emph{por qué}. Filas: \emph{CPU, RAM, disco, GPU, monitor, periféricos}. Para cada uno, escribe qué eliges (ejemplo: \emph{``CPU: Intel Core i5''}) y por qué con tus palabras (\emph{``Porque necesito que abra Word, Chrome y Minecraft con shaders sin pegarse''}).}
\milcpaso{3}{milcMagenta}{\textbf{Paso 3 --- Presupuesto estimado.} Debajo de la tabla, escribe \textbf{``Presupuesto aproximado:''} con un rango (\emph{``$2.000.000 a $2.800.000 pesos colombianos''}). Si no sabes precios, mira el rango sugerido por perfil y ajusta. La idea es que sepas que un computador no son $300.000 pero tampoco $20.000.000.}
\milcpaso{4}{milcMagenta}{\textbf{Paso 4 --- Software inicial.} Lista los \textbf{5 programas} que vas a usar más en este equipo (más el sistema operativo). Ejemplo: \emph{Sistema operativo: Windows 11. Programas: (1) Microsoft Word, (2) Chrome, (3) Spotify, (4) Roblox, (5) CapCut.}. Esta lista te recuerda que el hardware sin software no sirve.}
\milcpaso{5}{milcMagenta}{\textbf{Paso 5 --- Mantenimiento y ergonomía.} En la mitad inferior, escribe en 2-3 líneas \textbf{cómo lo vas a mantener} (ya sabes las acciones del S8) y \textbf{cómo lo vas a usar bien físicamente} (postura del S9). Esto cierra el círculo del periodo.}
\milcpaso{6}{milcMagenta}{\textbf{Paso 6 --- Justificación final + firma.} Al final escribe un párrafo de 3-4 líneas que diga: \emph{``Este computador no es el más caro ni el más barato. Es el que necesito para [tu perfil]. Si lo cuido bien, me va a durar X años. Lo elegí porque [razón principal].''}. Firma con nombre y fecha. Esta es \textbf{la prueba de que aprendiste a decidir con criterio}.}

\begin{drawbox}{milcMagenta}{Antes de mostrar tu ficha}
\checkbox Mi perfil está descrito con honestidad arriba. \\[1mm]
\checkbox La tabla de 6 componentes está completa con justificación. \\[1mm]
\checkbox El presupuesto tiene rango (no precio único exacto). \\[1mm]
\checkbox Lista de software inicial con SO + 5 programas. \\[1mm]
\checkbox Mantenimiento y ergonomía mencionados al menos en 2 líneas. \\[1mm]
\checkbox Justificación final firmada y fechada.
\end{drawbox}

\begin{softbox}{milcMagenta}{milcGris}{Plantilla de trabajo}
\textbf{Estructura.} \textit{Arriba:} título + mi perfil. \textit{Tabla:} 6 componentes (CPU/RAM/disco/GPU/monitor/periféricos) con elección y razón. \textit{Presupuesto:} rango en pesos. \textit{Software:} SO + 5 programas. \textit{Mantenimiento + ergonomía:} 2-3 líneas. \textit{Final:} justificación de 3-4 líneas + firma.
\end{softbox}

\begin{infoband}{milcMagenta}


\textbf{Lo que va al cuaderno (Actividad 3 · CREA).} Título: «Actividad 3 --- Mi computador ideal». \textbf{Formato:} ficha técnica completa con perfil, componentes, precio, software, mantenimiento y justificación. \textbf{Extensión:} 1 hoja entera. \textbf{Sabes que terminaste cuando} podrías ir a una tienda con esta hoja y comprar el equipo correcto, sin que te vendan algo distinto.
\end{infoband}

\puente{El producto es esa ficha completa con perfil, componentes, presupuesto, justificación y software.}

\milcestacion{milcMostaza}{Producto de la sesión}
\textbf{Ficha técnica del computador ideal · perfil + componentes + presupuesto + software · firmada}.
\begin{softbox}{milcMostaza}{milcGris}{Criterios de calidad}
(1) El perfil de usuario es honesto y específico. (2) Los 6 componentes están justificados por el perfil. (3) El presupuesto tiene rango realista. (4) El software inicial incluye SO + 5 programas. (5) La justificación final muestra criterio (no solo gusto).
\end{softbox}

\puente{Antes de salir, miras tu ficha: ¿coincide el perfil con los componentes? ¿La justificación tiene criterio?}

\milcestacion{milcVino}{Evaluación liberadora · cinco dimensiones}

\begin{softbox}{milcVino}{milcGris}{Sentido de la evaluación}
Evaluar no es solo revisar una nota. Es reconocer qué aprendí, qué cambió en mí, qué emociones manejé, cómo me relaciono con la ciudad y la comunidad, y qué le dejaré a quien viene detrás.
\end{softbox}

\textbf{Mi reflexión liberadora en cinco dimensiones.} Completa cada frase iniciada:

\small
\begin{tabularx}{\textwidth}{>{\bfseries\RaggedRight\arraybackslash}p{3.4cm}Y>{\RaggedRight\arraybackslash}p{4.6cm}}
\rowcolor{milcVino}\color{white}Dimensión & \color{white}\bfseries Mi respuesta (frame starter) & \color{white}\bfseries Algo que descubrí\\
1. Personal & Lo que más cambió en mí fue \linead{6.5cm} & \linead{4.2cm}\\
2. Emocional & La emoción más fuerte fue \linead{2.5cm} y la regulé \linead{3cm} & \linead{4.2cm}\\
3. Ciudadana & En democracia esto importa porque \linead{6cm} & \linead{4.2cm}\\
4. Local (Cartago) & En mi barrio sería útil para \linead{6.5cm} & \linead{4.2cm}\\
5. Intergeneracional & A mi abuela/abuelo le diría \linead{6.5cm} & \linead{4.2cm}\\
\end{tabularx}
\normalsize

\begin{infoband}{milcVino}
\textbf{Para la próxima sustentación.} Vuelve a esta tabla en seis meses. Verás que las respuestas cambian — no porque lo aprendido cambie, sino porque tú cambias. La evaluación liberadora es una conversación contigo a través del tiempo.
\end{infoband}


\puente{Cerramos con 3 voces que te ayudan a entender por qué decidir con criterio es libertad.}

\milcestacion{milcGradeDark}{Triángulo de pensamiento · cierre filosófico}

\begin{softbox}{milcVino}{milcGris}{Dussel · lente del nosotros}
\textit{Quien sabe qué necesita, escapa de la trampa del consumo. Quien no, compra lo que le venden, no lo que necesita.}

{\footnotesize\itshape\color{milcNegro!70}--- formulación de esta guía, al modo de Enrique Dussel. No es una cita textual.}

La publicidad de computadores te dice \emph{``compra el último modelo''}, \emph{``el más rápido''}, \emph{``el más cool''}. Pero TÚ ahora sabes que \textbf{no necesitas el último modelo}: necesitas el que se adapta a tu perfil. \emph{Saber esto es libertad}: no te dejas llevar por el marketing, decides con criterio. Esa es la diferencia entre consumir y elegir.

\textbf{¿Las cosas que quiero comprar las necesito, o me las quieren vender porque están de moda?}
\end{softbox}
\linead{\linewidth}\\[1mm]
\linead{\linewidth}

\begin{softbox}{milcMostaza}{milcGris}{Estoicismo · Marco Aurelio · cuidado interior}
\textit{Mide cada decisión por su propósito, no por el ruido a su alrededor.}

{\footnotesize\itshape\color{milcNegro!70}--- formulación de esta guía, al modo de Marco Aurelio. No es una cita textual.}

Cuando vas a comprar un computador (o cualquier cosa importante), \textbf{el ruido es enorme}: vendedores, publicidad, youtubers, amigos que dicen \emph{``compra el mismo que yo''}. El sabio se detiene y pregunta: \emph{``¿para qué lo necesito yo?''}. Si tu respuesta es clara, decides bien. Si no, te enredan.

\textbf{Mis últimas 3 compras importantes, ¿las decidí por mi propósito o por el ruido?}
\end{softbox}
\linead{\linewidth}\\[1mm]
\linead{\linewidth}

\begin{softbox}{milcTurquesa}{milcGris}{Floridi · lente de la infoesfera}
\textit{El equipo correcto te dura años. El equipo equivocado te genera el deseo de comprar otro en meses.}

{\footnotesize\itshape\color{milcNegro!70}--- formulación de esta guía, al modo de Luciano Floridi. No es una cita textual.}

Si compras un computador que NO se adapta a tu uso, vas a sentir frustración: \emph{``es muy lento''}, \emph{``no me aguanta lo que quiero hacer''}. La frustración te lleva a querer comprar otro. \textbf{El equipo correcto te dura 5-8 años sin necesidad de cambio}. Decidir bien hoy es ahorrar dinero, tiempo y recursos del planeta durante años.

\textbf{¿La compra que estoy pensando me va a durar años o me va a generar ganas de cambiar en meses?}
\end{softbox}
\linead{\linewidth}\\[1mm]
\linead{\linewidth}

\begin{infoband}{milcNegro}
\textbf{Nota para el docente.} Las tres formulaciones son del autor de la guía, no citas textuales de los pensadores: recogen su lente para pensar el tema, no sus palabras. Si vas a citarlos en clase, remite a la obra original.
\end{infoband}

\milcestacion{milcVino}{Mi autoevaluación · marca una casilla por fila}

{\small
\setlength{\extrarowheight}{2.2pt}
\begin{tabularx}{\textwidth}{>{\RaggedRight\arraybackslash\bfseries}p{3.0cm}YYY}
\rowcolor{milcVino}\color{white}\bfseries Criterio & \color{white}\bfseries Lo logré & \color{white}\bfseries En proceso & \color{white}\bfseries Necesito apoyo\\[.6mm]
Comprensión del tema & \milcopcion{Explico las ideas con mis palabras.} & \milcopcion{Reconozco las ideas, falta precisión.} & \milcopcion{Necesito repasar lo básico.}\\[.8mm]
\rowcolor{milcGris}Pensamiento computacional & \milcopcion{Usé los pilares al armar mi producto.} & \milcopcion{Usé algunos pilares.} & \milcopcion{Necesito releer la estación 2.}\\[.8mm]
Aplicación y producto & \milcopcion{Mi producto cumple los criterios.} & \milcopcion{Está armado, con ajustes pendientes.} & \milcopcion{Aún está en construcción.}\\[.8mm]
\rowcolor{milcGris}Reflexión 5 dimensiones & \milcopcion{Respondí cada una con honestidad.} & \milcopcion{Respondí algunas con detalle.} & \milcopcion{Mis respuestas son muy generales.}\\[.8mm]
Triángulo de pensamiento & \milcopcion{Conecté las 3 voces con mi trabajo.} & \milcopcion{Conecté al menos una.} & \milcopcion{No las relaciono con mi proceso.}\\[.8mm]
\end{tabularx}}

\vspace{3mm}
\textbf{Hoy entendí que:}\\[1mm]
\linead{\linewidth}\\[1mm]
\linead{\linewidth}

\textbf{Mi compromiso para la próxima sesión es:}\\[1mm]
\linead{\linewidth}\\[1mm]
\linead{\linewidth}

\begin{softbox}{milcMostaza}{milcGris}{Cita de despedida · Marco Aurelio}
\textit{``El obstáculo en el camino se vuelve el camino. Lo que estorba al camino se vuelve camino.''}\\[1mm]
Cada error de hoy, bien observado, se convierte en pieza del camino que pisarás mañana.
\end{softbox}

\small
\begin{tabularx}{\textwidth}{>{\bfseries}p{3.6cm}Y}
\rowcolor{milcGradeDark!12}Nombre completo: & \linead{10cm}\\
Fecha: & \linead{4cm} \quad Curso/grupo: \linead{4cm}\\
Mi firma: & \linead{9cm}\\
\end{tabularx}
\normalsize

\begin{infoband}{milcGradeDark}
\textbf{Para la próxima sesión.} Guarda tu ficha. Si algún día tus papás o tú van a comprar un computador, esta hoja es tu guía. La próxima sesión (S11) es la \textbf{cosecha del periodo}: te autoevalúas con rúbrica, vemos el examen integrador y miramos al periodo 3 que empieza (Procesador de texto e Internet). Hoy cerraste el ciclo del hardware. La próxima clase cierras formalmente el periodo.
\end{infoband}

\end{document}
